\section{Discussion}

The results in Section~3 support a coherent story, but the value of a graduation report lies not only in showing curves and tables. It lies in drawing the right conclusions from them and making the scope of those conclusions explicit. This section therefore turns the numerical evidence into a disciplined interpretation.

\subsection{Mechanism Interpretation}

\emph{Weighted phase variance} is central to this project. After the nominal constructive design is normalized by its amplitude weights, the coherent sum depends primarily on the spread of elementwise phase errors about their weighted mean. This framework accounts for all major observations more economically than the earlier same-sign perturbation narrative.

We first clarify why common-bias motion is nearly harmless. When all activated pinching antennas shift in roughly the same effective phase direction, the perturbation contributes mainly to a common phase term rather than to a relative spread. In the second-order expansion that common term enters the weighted mean and cancels from the variance. The common-bias curve remains almost exactly at unity, which corroborates the phase-variance interpretation; it is not an isolated numerical accident.

The same mechanism also explains the striking efficacy of the split-mode perturbation in the present symmetric configuration. The constructive placement renders the positive- and negative-side geometries nearly mirror-symmetric, and by symmetry the weighted mean of the sensitivity profile stays precisely at \(n_{\mathrm{eff}}\). Once \(\xi_n - \bar{\xi}_\alpha\) is centered about zero, the error pattern that maximizes the local phase spread naturally drives the right and left halves in opposite directions. The resulting adversary is conceptually simple yet physically meaningful: it achieves the largest relative phase separation across the aperture while respecting the error budget. The close agreement between the split mode and the box-best-found adversary before the first null confirms that this local mechanism is not merely descriptive—it is predictive in the pre-null regime.

The stochastic results provide further support. With independent perturbations the covariance energy sits on the diagonal, and the centered weighting in \eqref{eq:stochastic_predictor} therefore yields a large elementwise phase spread. Common-bias motion remains nearly invisible because its covariance has rank one in the all-ones direction and is largely annihilated by the centered operator. Spatially correlated Gaussian perturbations fall between these extremes, as their smooth structure induces partial common-mode and partial differential behavior. Thus the stochastic section is not a supplementary addendum; it constitutes a second, independent validation of the very same weighted phase-variance object that already rationalizes the deterministic split-mode behavior.

\subsection{Validity Condition Versus Empirical Region}

A central refinement introduced during this project is the explicit separation of consequences that follow from the deterministic small‑error validity condition and those that are merely empirical. The report preserves this distinction for both the lower bound and the worst‑case search.

Inside the small‑error region defined by \(k_0 \xi_{\max}\epsilon \le \pi/2\), the projection step employed in the lower‑bound derivation holds under the first‑order phase‑deviation model. Accordingly, the recorded threshold \(\epsilon/\lambda = 0.1713\) is not a fitted parameter; it is the direct output of the analytic condition on which the bound rests. That the lower bound coincides with the box‑best‑found curve to within approximately \(3.40 \times 10^{-6}\) over the entire interval constitutes strong evidence of tightness for the studied design. In this sense, the report’s strongest claim is neither the split‑mode account nor the baseline comparison, but the tight deterministic small‑error characterization valid within the stated region.

Beyond that threshold, the report deliberately shifts its terminology. The lower bound may remain close to the observed curve for some range, yet such agreement is empirical. Similarly, the continuous‑box local‑search routine can return values lower than the corner‑restricted adversary, but this does not certify a global continuous minimum. For this reason the term \emph{box‑best‑found} is used consistently. The language is conservative by necessity: the nonconvex bounded‑box problem is not globally solved here.

The near‑null point around \(\epsilon/\lambda = 0.175\) offers the cleanest illustration of why this distinction matters. Up to \(\epsilon/\lambda = 0.15\), the corner solution, the split‑mode construction, and the box‑best‑found local search all agree to high precision. A cursory reading might then conclude that the corners are always exact or that the split‑mode solution is exact. The \(0.175\) case shows why such a conclusion would be premature. There the local search reaches numerical zero while the corner and split‑mode candidates remain at about \(1.58 \times 10^{-4}\).
 The correct interpretation is that the pre-null local argument is very strong, but not globally universal. The report therefore treats the pre-null split-mode pattern as an empirical near-adversarial construction, not as a theorem.

\subsection{Comparative Design Takeaway}

The baseline section adds value, but only if its meaning is kept narrow. The constructive aligned placement should not be advertised as the best design in the same-aperture comparison, because it is not. The best-found nominal-gain comparator is stronger on nominal gain, on box-best-found robustness, and on iid mean robustness. This is a concrete result, not a minor caveat.

At the same time, the constructive aligned design clearly does outperform the naive uniform-aperture baseline, especially under box-best-found robustness. The pattern of the gains is informative. Improvements in nominal gain and iid mean are both around one percent, whereas the improvement in box-best-found robustness exceeds fourteen percent. This difference suggests that the constructive alignment rule is particularly good at controlling the harmful phase-spread direction identified by the split-mode analysis. In essence, the constructive design offers an improvement over the nominal placement that is more than marginal. Its phase sensitivity distribution appears to mitigate the box-best-found degradation relative to that observed in a uniformly spaced aperture of identical span.

For the purposes of the present report, the appropriate design-level conclusion can be stated as follows: the constructive phase-aligned feasible placement yields a meaningful gain over a naive same-aperture uniform baseline in the studied robustness problem, while falling short of the best placement tested overall. This claim is scientifically valuable precisely because its scope is limited. It communicates that the phase-aligned constructive design makes a non-vacuous contribution, and it simultaneously confirms that scope remains for further gains through robustness-aware or nominally optimized placement strategies.

\subsection{Engineering Implications of the Sensitivity Checks}

The secondary sweeps presented in Section~3.4 derive their value from translating the underlying local mechanism into engineering terms. The effective-index sweep is particularly instructive. Raising \(n_{\mathrm{eff}}\) does more than rescale a coordinate axis; it amplifies the guided-wave contribution within the PASS phase law itself. The associated decrease in the box-best-found normalized gain accordingly strengthens the interpretation that PASS robustness arises not solely from free-space geometry but from the joint geometry–waveguide phase accumulation that produces constructive alignment.

The frequency sweep offers a complementary insight. When a fixed absolute placement error is specified in meters rather than wavelengths, the effective perturbation grows more pronounced as the carrier frequency increases. Thus, a placement-accuracy budget that is adequate at \SI{6}{GHz} may become unacceptable at \SI{28}{GHz} or \SI{60}{GHz}, even when the same mechanical actuation hardware is employed. For PASS systems operating at higher frequencies, preserving an equivalent robustness margin would therefore require tighter position control, shorter effective apertures, or more extensive active calibration.

Taken together, the common-bias result and the sensitivity sweeps suggest a clear priority for future control strategies. Given that common drift is largely benign whereas differential motion is destructive, the most valuable correction effort may not be whole-aperture translation compensation. Instead, suppressing or calibrating the differential error modes that widen the centered phase spread would prove more effective. While this falls short of a closed-loop demonstration, it represents a concrete design insight derived from the current evidence.

These implications are subject to the same scope constraints as the remainder of the report. They constitute design-specific interpretations for the symmetric configuration examined, not universal PASS scaling laws. Their contribution lies in demonstrating that the main theoretical coefficients are more than abstract symbols—they also identify the conditions under which robustness pressure would be expected to intensify in future engineering realizations of the problem.


\subsection{Limitations and Unsupported Claims}

Several limitations warrant attention. First, the bounded-box continuous search yields only best-found solutions, not provably global optima. Although integrating exact corner enumeration with box-constrained local refinement markedly outperforms a corner-only strategy, the overall procedure remains a heuristic for the continuous formulation. This restriction is most consequential near the first null, where the objective landscape becomes sharply nonconvex.

Second, the split-mode construction is an empirically identified near-adversarial pattern rather than an exact minimizer statement. The weighted phase-variance derivation explains why the construction appears compelling, yet the reasoning is fundamentally local and confined to small-error regimes. Global optimality over the full box for arbitrary error magnitudes is not established.

Third, the stochastic predictor has been validated solely for the studied constructive design and only up to $\epsilon/\lambda = 0.10$. Within that interval the approximation error remains modest, but the theory does not claim to retain comparable accuracy under larger errors, substantially different placements, or arbitrary covariance structures.

Fourth, the comparative design evaluation is limited in scope. The baseline analysis does not sweep over $N$, $d$, or $\Delta_p$, and it includes no explicitly robustness-tuned alternative placement. Consequently, the observation that the constructive aligned placement outperforms a naive uniform baseline while underperforming the best-found nominal-gain comparator is a context-dependent finding rather than a universal design ranking.

Taken together, these constraints delineate the report’s scientific standing. The contribution is a meticulous robustness characterization of one constructive PASS design within a single symmetric configuration, grounded in reproducible simulations and a coherent local theory. It does not constitute a universal theorem of PASS robustness, nor does it survey the full PASS design space. The report’s strength is that it makes these boundaries explicit, rather than concealing them behind stronger rhetoric.

